\chapter{Conclusiones y trabajo futuro}
\label{chap:conclusiones-trabajo-futuro}

%==============================================================================
% 7.1. Conclusiones
% =============================================================================
\section{Conclusiones}
\label{sec:conclusiones}

El desarrollo de este proyecto ha permitido construir un sistema completo para disputar una partida física de Onitama entre un jugador humano y un agente de inteligencia artificial. La solución obtenida no se limita a implementar las reglas del juego o a seleccionar movimientos de forma automática, sino que conecta estas capacidades con la observación de un tablero real mediante visión por computador y con una interfaz que coordina la interacción durante la partida.

La principal aportación del trabajo se encuentra en la integración entre la dimensión física y la representación lógica del juego. Una imagen capturada por la cámara no constituye por sí sola un estado válido, por lo que ha sido necesario distinguir entre la observación visual y el estado confirmado de la partida. La estabilización de las observaciones y su comparación con las transiciones permitidas por el motor de reglas permiten mantener la coherencia entre ambas representaciones y evitar que lecturas incompletas o incorrectas modifiquen indebidamente la partida.

Esta separación de responsabilidades también ha resultado fundamental para la organización del sistema. El motor de juego, el agente de inteligencia artificial, el módulo de visión, la lógica de sincronización, el control de ejecución y la interfaz gráfica se han desarrollado como componentes diferenciados, pero coordinados mediante representaciones y operaciones comunes. Este enfoque ha facilitado su validación independiente y ha permitido abordar de forma progresiva la integración final.

Onitama ha resultado un caso de estudio adecuado para este propósito. Su escala permite representar y analizar la partida con recursos computacionales accesibles, mientras que la rotación de las cartas y las distintas condiciones de victoria mantienen una complejidad estratégica suficiente para justificar el desarrollo de un agente de búsqueda. Del mismo modo, el reducido número de piezas y cartas visibles hace viable la reconstrucción del estado mediante una cámara, sin eliminar los retos propios de una escena física real.

En conjunto, el trabajo muestra que es posible incorporar un agente inteligente a un juego de tablero físico sin sustituir sus componentes materiales por una interfaz completamente digital. El sistema desarrollado actúa como enlace entre la percepción visual, la aplicación de las reglas y la toma automática de decisiones, manteniendo al jugador humano como responsable de la manipulación del tablero. El resultado constituye una base funcional sobre la que pueden realizarse nuevas mejoras tanto en la robustez visual como en la autonomía y capacidad de decisión del sistema.


%==============================================================================
% 7.2. Cumplimiento de los objetivos
% =============================================================================
\section{Cumplimiento de los objetivos}
\label{sec:cumplimiento-objetivos}

El objetivo general definido al inicio de la memoria consistía en diseñar, implementar y evaluar un sistema inteligente que permitiera disputar una partida física de Onitama entre un jugador humano y un agente de inteligencia artificial. Tras el desarrollo realizado, este objetivo puede considerarse cumplido, ya que se ha construido una aplicación capaz de representar la partida, observar el tablero físico, validar la evolución del juego y coordinar la actuación del agente automático.

El cumplimiento de este objetivo general se concreta en los distintos objetivos específicos planteados:

\begin{itemize}
    \item \textbf{Representación digital de la partida.} Se ha definido una representación completa del estado de Onitama, incluyendo el tablero, las piezas, el jugador activo, las cartas disponibles para cada jugador y la carta lateral. Esta representación ha servido como base común para el motor de reglas, la inteligencia artificial y la sincronización con la observación visual.

    \item \textbf{Implementación del motor de reglas.} Se ha desarrollado un motor capaz de generar acciones legales, aplicar movimientos, resolver capturas, actualizar el intercambio de cartas y detectar las dos condiciones de victoria del juego. Este componente actúa como referencia lógica del sistema y garantiza que la evolución interna de la partida respete las reglas de Onitama.

    \item \textbf{Desarrollo del agente de inteligencia artificial.} Se ha implementado un agente basado en búsqueda adversarial, con funciones heurísticas de evaluación y distintas técnicas de optimización. Además, se han definido perfiles de dificultad que permiten ajustar el comportamiento del agente en función de la profundidad de búsqueda empleada.

    \item \textbf{Reconstrucción visual del estado físico.} Se ha construido un módulo de visión por computador capaz de procesar imágenes capturadas por una cámara, detectar las piezas sobre el tablero y clasificar las cartas visibles. A partir de esta información, el sistema puede generar una observación estructurada de la escena física.

    \item \textbf{Validación de los cambios observados.} Se ha incorporado una capa de sincronización que no acepta directamente cualquier observación visual, sino que la estabiliza y la compara con las transiciones legales generadas por el motor. De este modo, el sistema puede distinguir entre estados válidos, observaciones sin cambio y configuraciones incoherentes con la evolución de la partida.

    \item \textbf{Coordinación de los turnos de juego.} Se ha desarrollado una lógica de sesión que coordina el turno del jugador humano, la selección de la jugada de la inteligencia artificial y la confirmación de su ejecución física sobre el tablero. Esta coordinación permite que el agente participe en una partida real sin que el juego se reduzca a una simulación puramente digital.

    \item \textbf{Desarrollo de la interfaz gráfica.} Se ha implementado una aplicación de escritorio que permite configurar la partida, realizar las calibraciones necesarias, visualizar el estado del juego, mostrar las cartas y guiar al usuario durante la interacción con el agente automático.

    \item \textbf{Evaluación experimental.} Se han realizado pruebas sobre los principales componentes del sistema. En particular, se ha evaluado el comportamiento del agente mediante torneos y benchmarks de búsqueda, se han entrenado y validado los modelos de visión por computador y se ha comprobado el funcionamiento conjunto de la aplicación en condiciones reales de uso.
\end{itemize}

En conjunto, los objetivos planteados se han alcanzado dentro del alcance definido para el proyecto. El sistema resultante permite jugar físicamente a Onitama con asistencia de inteligencia artificial y visión por computador, aunque su funcionamiento sigue condicionado por las restricciones propias del montaje físico, la cámara y los modelos visuales empleados.


%==============================================================================
% 7.3. Principales resultados
% =============================================================================
\section{Principales resultados}
\label{sec:principales-resultados}

La evaluación experimental permitió comparar las distintas alternativas desarrolladas y justificar la configuración final del sistema. Los resultados más relevantes se resumen a continuación.

\paragraph{Comparación de las funciones heurísticas.}
En el agente de inteligencia artificial, la función heurística \texttt{v3} obtuvo el mejor comportamiento en los cuatro torneos realizados. Su tasa de puntuación se situó alrededor del 65\,\% en profundidad 3 y entre el 72\,\% y el 76\,\% en profundidad 4, superando tanto a \texttt{v1} como a \texttt{v2} en los enfrentamientos directos. Estos resultados respaldan el uso de una evaluación dependiente de la fase de la partida, capaz de adaptar la importancia de los distintos criterios estratégicos conforme evoluciona el juego.

\paragraph{Impacto de las optimizaciones de búsqueda.}
El análisis realizado mostró que no todas las optimizaciones implementadas aportaban el mismo beneficio en el escenario estudiado. La ordenación de movimientos fue la técnica con un efecto más claro, ya que su desactivación incrementó progresivamente el tiempo y el número de nodos explorados al aumentar la profundidad. La tabla de transposición presentó un comportamiento menos estable, mientras que la profundización iterativa y las ventanas de aspiración añadieron coste en las búsquedas a profundidad fija utilizadas. La configuración final, que conserva la ordenación de movimientos y desactiva estas tres técnicas, redujo el tiempo total respecto a la configuración completa en un 18.4\,\%, un 19.4\,\% y un 31.2\,\% para las profundidades 3, 4 y 5, respectivamente.

\paragraph{Efecto de la búsqueda quiescente.}
La búsqueda quiescente produjo una mejora clara en la calidad del agente. En los torneos comparativos, la configuración sin esta extensión obtuvo una tasa de puntuación próxima al 33\,\%, mientras que la profundidad quiescente 2 alcanzó valores cercanos al 62\,\% en los dos bloques de semillas evaluados. El aumento del tiempo por decisión fue reducido dentro del contexto de una partida interactiva, por lo que esta configuración se mantuvo en los perfiles definitivos.

\paragraph{Perfiles de dificultad.}
Los torneos entre niveles confirmaron una separación coherente entre los perfiles definidos. El perfil fácil obtuvo una tasa de puntuación del 12.38\,\%, el perfil medio alcanzó el 49.88\,\% y el perfil difícil llegó al 87.75\,\%. Los tiempos medios por decisión permanecieron por debajo de una décima de segundo incluso en el nivel difícil. La combinación de la heurística \texttt{v3}, una profundidad principal de 5 y una profundidad quiescente de 2 permitió así alcanzar un nivel de juego elevado sin perjudicar la fluidez de la partida. En las partidas manuales realizadas durante el desarrollo, este perfil ofreció además un reto considerable frente a jugadores familiarizados con Onitama.

% Los apartados sobre los modelos de visión y la valoración conjunta
% se completarán cuando finalicen las pruebas en condiciones reales
% y la validación funcional del sistema integrado.
Faltan los dos últimos apartados!!

%==============================================================================
% 7.4. Limitaciones
% =============================================================================
\section{Limitaciones}
\label{sec:limitaciones}

Aunque el sistema desarrollado cumple los objetivos planteados, su funcionamiento está condicionado por varios factores que deben tenerse en cuenta al interpretar los resultados y valorar su posible uso en otros entornos.

\paragraph{Condiciones de captura y generalización visual.}
El rendimiento del módulo de visión depende de factores como la iluminación, el enfoque, los reflejos, las oclusiones y la visibilidad de las piezas y cartas. Además, los modelos se han entrenado con los componentes físicos y las condiciones de captura empleadas durante el desarrollo. Por ello, cambios significativos en el tablero, las piezas, las cartas o el entorno podrían reducir la fiabilidad y requerir nuevos datos de entrenamiento o ajustes adicionales.

\paragraph{Calibración y montaje físico.}
El sistema requiere una calibración previa del tablero y de las regiones correspondientes a las cartas. Durante la partida se asume que la cámara permanece fija y que la disposición general del montaje no cambia. Si se modifica el encuadre, la posición del tablero o la ubicación de las cartas, es necesario repetir la calibración. Esta dependencia limita la capacidad de utilizar la aplicación de forma inmediata en configuraciones físicas diferentes.

\paragraph{Sincronización de la partida.}
La lógica de sincronización permite estabilizar las observaciones y comprobar si los cambios detectados corresponden a transiciones legales desde el último estado confirmado. Sin embargo, no contempla la recuperación automática ante desincronizaciones complejas, como varios movimientos no observados o modificaciones físicas que no puedan relacionarse con una única acción legal. En estas situaciones, el usuario debe corregir manualmente la disposición o reiniciar la sesión.

\paragraph{Evaluación del agente.}
La calidad del agente se ha analizado principalmente mediante enfrentamientos automáticos entre distintas heurísticas, configuraciones de búsqueda y perfiles de dificultad. Este enfoque permite establecer comparaciones relativas entre las alternativas desarrolladas, pero no determina un nivel absoluto frente a jugadores humanos ni frente a otros motores de Onitama. Asimismo, el coste computacional de profundidades superiores dificultó la ejecución de torneos con un número de partidas comparable al utilizado en las configuraciones finalmente seleccionadas.

\paragraph{Autonomía del sistema.}
La aplicación calcula y muestra las acciones de la inteligencia artificial, pero no incluye mecanismos capaces de mover físicamente las piezas y cartas. Por tanto, el jugador humano debe ejecutar sobre el tablero tanto sus propias jugadas como las indicadas por el agente. El sistema actúa como observador, validador y oponente lógico, pero no como un jugador físico completamente autónomo.

%==============================================================================
% 7.5. Trabajo futuro
% =============================================================================
\section{Trabajo futuro}
\label{sec:trabajo-futuro}

Las limitaciones identificadas permiten plantear distintas líneas de continuación orientadas a mejorar la robustez, la flexibilidad y la autonomía del sistema desarrollado.

\paragraph{Mejora del módulo de visión.}
Una primera línea de trabajo consiste en ampliar los conjuntos de datos con nuevas condiciones de iluminación, cámaras, orientaciones y situaciones de oclusión. Esto permitiría evaluar mejor la capacidad de generalización de los modelos y reducir su dependencia del montaje utilizado durante el desarrollo. También podría incorporarse información temporal entre frames o una detección explícita de manos y oclusiones para evitar procesar escenas mientras el jugador está manipulando los componentes.

\paragraph{Calibración más flexible.}
El sistema podría reducir su dependencia de una configuración fija mediante mecanismos automáticos de localización del tablero y de las cartas. El uso de referencias visuales, marcadores o técnicas de seguimiento permitiría actualizar la transformación geométrica cuando se produjeran pequeños desplazamientos de la cámara o del tablero, evitando repetir manualmente toda la calibración.

\paragraph{Recuperación de la sincronización.}
La lógica actual puede detectar observaciones incoherentes, pero no reconstruir automáticamente una partida cuando se pierde la correspondencia con el último estado confirmado. Como mejora, podrían analizarse secuencias de varias acciones legales para encontrar una posible ruta hasta la posición observada. También resultaría útil permitir la corrección manual del estado desde la interfaz y proporcionar indicaciones más precisas sobre las piezas o cartas responsables de la discrepancia.

\paragraph{Ampliación del agente de inteligencia artificial.}
El agente podría extenderse mediante nuevas funciones de evaluación, ajuste automático de pesos o estrategias de búsqueda basadas en límites de tiempo por movimiento. También sería posible estudiar enfoques alternativos, como la búsqueda Monte Carlo o métodos de aprendizaje por refuerzo. Finalmente, una evaluación con jugadores de distintos niveles o frente a otras implementaciones de Onitama permitiría situar con mayor precisión la capacidad de juego del agente.

\paragraph{Mejoras en la interacción y autonomía.}
La interfaz podría incorporar un historial completo de movimientos, reproducción de partidas, guardado y recuperación de sesiones y herramientas más avanzadas de diagnóstico visual. Como ampliación de mayor alcance, podría integrarse un mecanismo robótico capaz de ejecutar físicamente las jugadas seleccionadas por la inteligencia artificial. De este modo, el sistema dejaría de depender del jugador humano para mover las piezas y se convertiría en un oponente físico completamente autónomo.